\chapter{One request, end to end}
\chapterintro{Before building a server, build a precise story of what a server participates in. A URL typed into a browser crosses many boundaries; HTTP is only one layer of that journey.}

\begin{objectives}
\begin{itemize}
  \item Identify the client, origin server, resolver, routers, and intermediaries in a web request.
  \item Separate name resolution, connection establishment, security, HTTP exchange, and application work.
  \item Distinguish a protocol from an implementation and a resource from its representation.
  \item Locate failures by phase instead of calling every problem ``the network.''
\end{itemize}
\end{objectives}

\section{Begin with a URL}

Imagine entering \url{https://api.example.test/tasks/17?view=compact} into a
client. The string is an instruction, not a connection. It identifies a scheme
(\code{https}), an authority (\code{api.example.test}), a path
(\code{/tasks/17}), and a query (\code{view=compact}). The client must turn that
description into communication with some process.

The scheme determines important defaults and rules. HTTPS normally implies a
secure connection and a default port of 443. The hostname is meaningful to
humans and to the Domain Name System, but the Internet Protocol forwards
packets using numeric addresses. The path and query are carried inside the HTTP
request; routers in the network do not use them to choose a physical route.

\begin{mentalmodel}
A URL answers \emph{what origin and resource do I want?}; DNS helps answer
\emph{which address can reach that origin?}; TCP creates an ordered byte
stream to a port; TLS protects that stream; HTTP gives the bytes message
semantics; application code decides what the resource means.
\end{mentalmodel}

\section{The phases}

A typical first request proceeds through the following phases. Real clients
reuse cached DNS answers, existing connections, and cached responses, so some
phases can disappear on later requests.

\begin{enumerate}
  \item The client parses and validates the URL.
  \item It consults local configuration and caches, then asks a DNS resolver for
        one or more IP addresses.
  \item The operating system selects a route and a source address.
  \item The client asks the operating system to open a TCP connection to the
        selected address and port.
  \item For HTTPS, client and server perform a TLS handshake. The client checks
        that the certificate is valid for the requested hostname.
  \item The client serializes an HTTP request into bytes and writes them to the
        connection.
  \item The server operating system accepts the connection and makes received
        bytes available to a process through a socket.
  \item Server code reads, frames, parses, validates, routes, and handles the
        request. It might query a database or another service.
  \item The server serializes an HTTP response and writes it to the connection.
  \item The client parses the response, interprets its status and fields, and
        uses the representation. A browser might then request more resources.
\end{enumerate}

\begin{figure}[htbp]
\centering
\begin{tikzpicture}[
  node distance=8mm and 8mm,
  box/.style={draw=Rule,fill=PaperTint,rounded corners=1mm,minimum height=9mm,
              text width=25mm,align=center,font=\small\sffamily},
  arrow/.style={-{Latex[length=2mm]},draw=Accent}
]
\node[box] (url) {URL and client};
\node[box,right=of url] (dns) {DNS resolver};
\node[box,right=of dns] (tcp) {TCP connection};
\node[box,right=of tcp] (tls) {TLS channel};
\node[box,below=of tls] (http) {HTTP parser};
\node[box,left=of http] (router) {Application router};
\node[box,left=of router] (db) {Resource and storage};
\draw[arrow] (url) -- node[above,font=\scriptsize]{name} (dns);
\draw[arrow] (dns) -- node[above,font=\scriptsize]{address} (tcp);
\draw[arrow] (tcp) -- node[above,font=\scriptsize]{bytes} (tls);
\draw[arrow] (tls) -- (http);
\draw[arrow] (http) -- node[above,font=\scriptsize]{request} (router);
\draw[arrow] (router) -- node[above,font=\scriptsize]{operation} (db);
\draw[arrow,dashed] (db.south) |- node[below,font=\scriptsize]{response travels back} (url.south);
\end{tikzpicture}
\caption{The conceptual phases of a first HTTPS request. Each arrow hides a
separate interface and possible failure.}
\label{fig:request-journey}
\end{figure}

\section{Client, server, and resource}

\term{Client} and \term{server} describe roles in one interaction. A process
can be a server while receiving an API request and then become a client when it
calls a database proxy or another service. The machine itself is not inherently
the server; the listening process is performing the server role.

An HTTP \term{resource} is an abstraction identified by a URI. A task is not
the JSON text sent across the wire. JSON is one \term{representation} of that
task at a particular time. The same resource might have an HTML representation,
a compact JSON representation, or no transferable representation after a
successful deletion. This distinction later explains media types, content
negotiation, validators, and status codes.

\section{Where the server begins and ends}

Our Python code will not implement Ethernet, IP, or TCP. The operating system
does that work and exposes a socket interface. ``From scratch'' therefore
means from the lowest portable application interface: raw TCP sockets. We will
manually define the HTTP boundary, parse requests, serialize responses, route
resources, compose middleware, control concurrency, and connect the application
to SQLite.

This boundary is important. Calling \code{recv} does not retrieve ``one HTTP
request.'' It retrieves some bytes currently available from a TCP stream. The
HTTP implementation must decide whether those bytes contain a complete header,
how long the body is, and whether bytes after that body belong to a later
request.

\section{Failure belongs to a phase}

Precise phase language makes debugging faster.

\begin{center}
\begin{tabularx}{\textwidth}{@{}>{\bfseries}p{28mm} X X@{}}
\toprule
Phase & Typical symptom & First question \\
\midrule
DNS & Name cannot be resolved & Did the resolver return an address? \\
Routing & Destination appears unreachable & Is there a route to that address? \\
TCP & Connection refused or times out & Is a process listening on the chosen address and port? \\
TLS & Certificate or handshake failure & Does identity verification match the hostname? \\
HTTP framing & Client waits, closes, or reports a malformed response & Are lines, lengths, and boundaries correct? \\
Application & A valid 4xx or 5xx response & Which resource rule or dependency failed? \\
\bottomrule
\end{tabularx}
\end{center}

\begin{pitfall}
An HTTP \status{404 Not Found} proves that an HTTP response came back. DNS,
routing, TCP, and enough HTTP parsing already worked. Treating it as a generic
connectivity failure sends debugging in the wrong direction.
\end{pitfall}

\section{Latency is accumulated waiting}

The total time observed by a user includes name resolution, connection and TLS
round trips, queueing, application work, database time, response transfer, and
client processing. A server handler that takes 20 milliseconds does not imply a
20-millisecond user experience. Conversely, connection reuse can remove several
round trips even when application code is unchanged.

This is the first systems lesson of the book: behavior emerges from a chain of
components. Measurements must state where they begin and end.

\summarytitle

A URL initiates several distinct protocols and interfaces. DNS supplies an
address; IP forwards packets; TCP exposes an ordered byte stream; TLS can protect
it; HTTP defines messages; application code gives resources meaning. Our server
will begin at the socket boundary and implement the layers above it.

\reviewtitle
\begin{enumerate}
  \item Why does a router not need the path \code{/tasks/17} to forward an IP packet?
  \item Which phases have already succeeded when a client receives a valid 404 response?
  \item Why is a JSON document a representation rather than the resource itself?
  \item In what situation can one process be both an HTTP server and an HTTP client?
\end{enumerate}

\begin{exercise}[Draw the journey]
Draw the request journey for a local URL such as
\url{http://127.0.0.1:8080/health}. Cross out the phases that do not occur and
explain why. Then redraw it for an HTTPS API reached through a reverse proxy.
\end{exercise}

\begin{exercise}[Classify failures]
Classify each symptom by its earliest likely phase: ``name not known,''
``connection refused,'' ``certificate hostname mismatch,'' ``400 Bad Request,''
and ``database busy.'' For each, name one observation that would confirm the
classification.
\end{exercise}

